\chapter{Fundamentos Matemáticos para Ecossistemas Cognitivos}

\begin{quote}
\textit{``A Matemática é a linguagem com a qual Deus escreveu o universo.''} --- Galileu Galilei
\end{quote}

\section{Álgebra Linear: Embeddings, Atenção e Representações Vetoriais}

\subsection{Espaços Vetoriais e o Significado Geométrico dos Embeddings}

Um \textbf{embedding} é uma representação vetorial densa que mapeia entidades discretas (palavras, agentes, conceitos) para um espaço contínuo $\mathbb{R}^d$. A mágica dos embeddings está nas \textbf{relações geométricas} que emergem: $v_{\text{rei}} - v_{\text{homem}} + v_{\text{mulher}} \approx v_{\text{rainha}}$ (Mikolov et al., 2013) \cite{mikolov2013efficient}.

Dada uma matriz de embeddings $E \in \mathbb{R}^{n \times d}$ com $n$ tokens e dimensão $d$:

\begin{equation}
E = \begin{bmatrix} \mathbf{e}_1 \\ \mathbf{e}_2 \\ \vdots \\ \mathbf{e}_n \end{bmatrix}, \quad \mathbf{e}_i \in \mathbb{R}^d
\label{eq:embedding_matrix}
\end{equation}

A similaridade entre dois tokens é medida pelo \textbf{coseno}:

\begin{equation}
\text{sim}(\mathbf{e}_i, \mathbf{e}_j) = \frac{\mathbf{e}_i \cdot \mathbf{e}_j}{\|\mathbf{e}_i\| \|\mathbf{e}_j\|}
\label{eq:cosine_sim}
\end{equation}

\subsection{O Mecanismo de Atenção como Álgebra Linear}

O \textbf{Scaled Dot-Product Attention} (Vaswani et al., 2017) \cite{vaswani2017attention} é uma composição de operações matriciais elementares:

\begin{equation}
\text{Attention}(Q, K, V) = \text{softmax}\left(\frac{QK^T}{\sqrt{d_k}}\right) V
\label{eq:attention}
\end{equation}

Onde $Q, K, V \in \mathbb{R}^{n \times d}$ são as matrizes de Query, Key e Value. O fator $\frac{1}{\sqrt{d_k}}$ evita que o produto escalar exploda para dimensões altas, mantendo a variância em 1.

\subsection{Decomposição em Valores Singulares (SVD) e LoRA}

A técnica \textbf{Low-Rank Adaptation} (LoRA; Hu et al., 2022) \cite{hu2022lora} explora o fato de que atualizações de pesos em fine-tuning têm baixo rank:

\begin{equation}
\Delta W = BA, \quad B \in \mathbb{R}^{d \times r}, A \in \mathbb{R}^{r \times k}, r \ll \min(d,k)
\label{eq:lora}
\end{equation}

Para uma matriz de pesos $W_0 \in \mathbb{R}^{d \times k}$, a atualização LoRA é $W = W_0 + \Delta W = W_0 + BA$, reduzindo o número de parâmetros treináveis de $d \times k$ para $r(d+k)$ — tipicamente 0.01\% do original.

\subsection{PRÁXIS 3.1 — Atenção e LoRA em PyTorch}

\begin{lstlisting}[language=Python, caption={Implementação de Self-Attention + LoRA}]
import torch
import torch.nn as nn
import math

class ScaledDotProductAttention(nn.Module):
    def __init__(self, d_k):
        super().__init__()
        self.d_k = d_k
    
    def forward(self, Q, K, V, mask=None):
        scores = torch.matmul(Q, K.transpose(-2, -1)) / math.sqrt(self.d_k)
        if mask is not None:
            scores = scores.masked_fill(mask == 0, -1e9)
        attn = torch.softmax(scores, dim=-1)
        return torch.matmul(attn, V), attn
\end{lstlisting}

\section{Cálculo e Otimização para Aprendizado de Máquina}

\subsection{Do Gradiente Descendente ao AdamW}

O \textbf{Gradient Descent} é o cavalo de batalha do aprendizado de máquina. A regra de atualização é:

\begin{equation}
\theta_{t+1} = \theta_t - \eta \nabla_\theta \mathcal{L}(\theta_t)
\label{eq:gd}
\end{equation}

\textbf{Adam} (Kingma \& Ba, 2015) \cite{kingma2015adam} combina momento e adaptação por parâmetro:

\begin{align}
m_t &= \beta_1 m_{t-1} + (1-\beta_1) g_t \\
v_t &= \beta_2 v_{t-1} + (1-\beta_2) g_t^2 \\
\hat{m}_t &= \frac{m_t}{1-\beta_1^t}, \quad \hat{v}_t = \frac{v_t}{1-\beta_2^t} \\
\theta_{t+1} &= \theta_t - \eta \frac{\hat{m}_t}{\sqrt{\hat{v}_t} + \epsilon}
\label{eq:adam}
\end{align}

\textbf{AdamW} (Loshchilov \& Hutter, 2019) \cite{loshchilov2019adamw} corrige o weight decay do Adam, separando-o da taxa de aprendizado adaptativa, melhorando a generalização.

\subsection{PRÁXIS 3.2 — AdamW do Zero}

\begin{lstlisting}[language=Python, caption={Implementação de AdamW}]
class AdamW:
    def __init__(self, params, lr=1e-3, betas=(0.9, 0.999), eps=1e-8, weight_decay=0.01):
        self.params = list(params)
        self.lr, self.betas, self.eps, self.wd = lr, betas, eps, weight_decay
        self.m = [torch.zeros_like(p) for p in self.params]
        self.v = [torch.zeros_like(p) for p in self.params]
        self.t = 0
    
    def step(self):
        self.t += 1
        for i, p in enumerate(self.params):
            if p.grad is None: continue
            g = p.grad
            self.m[i] = self.betas[0] * self.m[i] + (1 - self.betas[0]) * g
            self.v[i] = self.betas[1] * self.v[i] + (1 - self.betas[1]) * g * g
            m_hat = self.m[i] / (1 - self.betas[0] ** self.t)
            v_hat = self.v[i] / (1 - self.betas[1] ** self.t)
            p.data -= self.lr * (m_hat / (torch.sqrt(v_hat) + self.eps) + self.wd * p.data)
\end{lstlisting}

\section{Teoria dos Grafos para Arquiteturas Multiagentes}

\subsection{Agentes como Vértices, Interações como Arestas}

Um ecossistema de 128 agentes pode ser modelado como um grafo dirigido $G = (V, E)$ onde vértices são agentes e arestas representam interações (chamadas de função, trocas no Blackboard, broadcasts no MetaBus).

\begin{equation}
A \in \{0,1\}^{n \times n}, \quad A_{ij} = \begin{cases} 1 & \text{se agente } i \rightarrow j \\ 0 & \text{caso contrário} \end{cases}
\label{eq:adjacency}
\end{equation}

\subsection{Métricas de Centralidade para Roteamento de Atenção}

O \textbf{Attention Router} do OpenCode implementa um análogo ao \textbf{PageRank} (Page et al., 1999) \cite{page1999pagerank}:

\begin{equation}
\text{TrustRank}(a_i) = \alpha \sum_{j \in \text{neighbors}(i)} \frac{\text{TrustRank}(a_j)}{\text{outdegree}(a_j)} + (1-\alpha) \cdot \text{TrustScore}(a_i)
\label{eq:trustrank}
\end{equation}

Onde $\text{TrustScore}(a_i)$ é a confiança intrínseca do agente $a_i$ (Capítulo 9) e $\alpha$ controla o peso da estrutura do grafo versus mérito individual.

\begin{thebibliography}{99}
\bibitem{mikolov2013efficient} Mikolov, T., et al. (2013). Efficient Estimation of Word Representations in Vector Space. \textit{arXiv:1301.3781}. \doiurl{10.48550/arXiv.1301.3781}
\bibitem{vaswani2017attention} Vaswani, A., et al. (2017). Attention Is All You Need. \textit{NeurIPS 2017}. \doiurl{10.48550/arXiv.1706.03762}
\bibitem{hu2022lora} Hu, E. J., et al. (2022). LoRA: Low-Rank Adaptation of Large Language Models. \textit{ICLR 2022}. \doiurl{10.48550/arXiv.2106.09685}
\bibitem{kingma2015adam} Kingma, D. P., \& Ba, J. (2015). Adam: A Method for Stochastic Optimization. \textit{ICLR 2015}. \doiurl{10.48550/arXiv.1412.6980}
\bibitem{loshchilov2019adamw} Loshchilov, I., \& Hutter, F. (2019). Decoupled Weight Decay Regularization. \textit{ICLR 2019}. \doiurl{10.48550/arXiv.1711.05101}
\bibitem{page1999pagerank} Page, L., Brin, S., Motwani, R., \& Winograd, T. (1999). The PageRank Citation Ranking. \textit{Stanford InfoLab}.
\bibitem{strang2016algebra} Strang, G. (2016). \textit{Introduction to Linear Algebra} (5th ed.). Wellesley-Cambridge Press.
\bibitem{bengio2003neural} Bengio, Y., et al. (2003). A Neural Probabilistic Language Model. \textit{JMLR}, 3, 1137--1155. \doiurl{10.5555/944919.944966}
\bibitem{barabasi2016network} Barabási, A.-L. (2016). \textit{Network Science}. Cambridge University Press.
\bibitem{newman2018networks} Newman, M. (2018). \textit{Networks} (2nd ed.). Oxford University Press.
\end{thebibliography}
